%% -----------------------------------------------------------------
%% 北京邮电大学研究生学位论文 LaTeX 模板 - 主文件
%% 遵循标准：北京邮电大学研究生院《学位论文写作与制作规范》（2025年1月15日修订）
%% 编译引擎：XeLaTeX + BibTeX
%% -----------------------------------------------------------------

%% =================================================================
%% 1. 文档类设置（请根据个人实际情况选择对应的文档类参数）
%% =================================================================
%% 学术学位硕士（电子版，默认）：
\documentclass[master,phil,fulltime,electronic]{buptthesis}

%% 其它可选类型（取消相应注释即可使用）：
%% \documentclass[master,phil,fulltime,print]{buptthesis}       % 学硕双面打印版
%% \documentclass[master,eng,fulltime,electronic]{buptthesis}    % 专硕电子版
%% \documentclass[master,eng,fulltime,print]{buptthesis}         % 专硕双面打印版
%% \documentclass[phd,phil,fulltime,electronic]{buptthesis}       % 学博电子版
%% \documentclass[phd,phil,fulltime,print]{buptthesis}            % 学博双面打印版
%% \documentclass[phd,eng,fulltime,electronic]{buptthesis}        % 专博电子版
%% \documentclass[phd,eng,fulltime,print]{buptthesis}             % 专博双面打印版

%% 选项说明：
%% master / phd        : 硕士学位 / 博士学位
%% phil / eng          : 学术学位 / 专业学位
%% fulltime / parttime : 全日制 / 非全日制
%% electronic / print  : 电子版（无多余空白页，适于送审与系统格式检测）/ 打印版（双面印刷，偶数页自动留白）

%% =================================================================
%% 2. 用户宏包与自定义配置（可选）
%% =================================================================
\usepackage{siunitx}      % 物理量与国际单位支持
\usepackage{subcaption}   % 子图排版支持

%% =================================================================
%% 3. 论文元数据填写（请按实际情况填写）
%% =================================================================

%% 中文扉页信息
\BUPTthesiscntitlepage{%
  confidential = {公开},                  % 密级：仅限涉密学生改写，默认“公开”
  title = {基于深度学习的高性能协同通信网络优化方法}, % 论文题目（不超过36个汉字）
  studentid = {2022110123},               % 10位学号
  name = {张\quad 三},                    % 作者姓名
  major = {信息与通信工程},               % 学科专业（一级/二级学科或专业学位领域）
  trainingmode = {全日制},                % 培养方式：全日制 / 非全日制
  supervisor = {李\quad 四\quad 教授},    % 导师姓名与职称
  institute = {信息与通信工程学院}         % 所在学院
}

%% 日期配置
\BUPTdateOptions{%
  pageYear = {2025},                     % 扉页年份
  pageMonth = {5},                       % 扉页月份
  pageDay = {30},                        % 扉页日期
  dissertationYear = {2025},             % 答辩年份
  dissertationMonth = {5},               % 答辩月份
  dissertationDay = {30}                 % 答辩日期
}

%% 英文扉页信息
\BUPTthesisentitlepage{%
  confidentialEN = {Public},
  titleEN = {Research on High-Performance Collaborative Communication Optimization via Deep Learning},
  studentidEN = {2022110123},
  nameEN = {San Zhang},
  majorEN = {Information and Communication Engineering},
  supervisorEN = {Prof. Si Li},
  instituteEN = {School of Information and Communication Engineering}
}

%% 答辩委员会名单（送审前版本可留空；答辩后终版如实填写）
\BUPTdissertationcommitteepage{%
  chairName = {王\quad 五},
  chairTitle = {教授},
  chairOrgan = {北京邮电大学},
  memberAName = {赵\quad 六},
  memberATitle = {教授},
  memberAOrgan = {北京邮电大学},
  memberBName = {孙\quad 七},
  memberBTitle = {研究员},
  memberBOrgan = {中国科学院计算技术研究所},
  memberCName = {周\quad 八},
  memberCTitle = {副教授},
  memberCOrgan = {北京邮电大学},
  memberDName = {吴\quad 九},
  memberDTitle = {教授},
  memberDOrgan = {清华大学},
  secretaryName = {郑\quad 十},
  secretaryTitle = {副教授},
  secretaryOrgan = {北京邮电大学}
}

%% 独创性声明扫描替换页（可选：如果已有签字盖章的整页扫描版PDF，取消下行注释指定文件路径即可）
%% \statementpage{resources/signed_statement.pdf}

%% =================================================================
%% 4. 论文正文排版架构
%% =================================================================
\begin{document}

%% ---------------- 扉页与版权声明部分 ----------------
\makechinesetitle      % 中文扉页（电子版即为第1页）
\makeenglishtitle      % 英文扉页
\makecommitteelist     % 答辩委员会名单
\makestatement         % 独创性声明与论文使用授权说明

%% ---------------- 前言部分（大写罗马数字页码） ----------------
\frontmatter
%% -----------------------------------------------------------------
%% 北京邮电大学研究生学位论文 - 中文摘要
%% 格式要求（2025年1月15日规范）：
%% 1. “摘要”标题黑体三号字，居中对齐，首行无缩进，段前0行，段后1行，单倍行距，中间空1个全角空格
%% 2. 摘要内容用四号宋体，两端对齐，首行缩进2字符，段前0行，段后0行，固定值20磅
%% 3. “关键词：”黑体四号字，其后关键词为宋体四号，两端对齐，首行缩进2字符，段前1行，段后0行，固定值20磅
%% 4. 关键词3～5个，每一关键词之间用全角分号隔开（；），最后一个关键词后不打标点符号
%% -----------------------------------------------------------------

\begin{cnabstract}
摘要应简明表达学位论文的内容要点，体现研究工作的核心思想。重点说明本项科研的目的和意义、研究方法、研究成果、结论，注意突出具有创新性的成果和新见解的部分。文字应精练、确切，中文摘要字数一般在千字左右。

本论文针对当前信息通信与智能计算领域的关键科学与工程问题，深入开展了理论建模、算法设计与系统验证研究。首先，系统梳理了国内外相关领域的研究现状与技术瓶颈；其次，提出了一种高鲁棒性、低时延的协同处理机制与自适应优化框架；接着，通过严谨的数学推导论证了所提方法在收敛速度与系统容量上的理论界限；最后，基于大规模真实数据集与原型平台开展了充分的对比实验，全面评估了系统在不同网络负载与动态环境下的运行表现。

实验结果表明，所提出的创新方案能够显著降低计算开销，提升吞吐量与决策精度，满足新一代信息处理系统对高可靠性与超低时延的严苛指标要求。本文的工作为该领域的进一步理论演进与工业应用提供了坚实的理论依据与参考方案。

\cnkeywords{北京邮电大学；研究生学位论文；\LaTeX{}模板；格式规范；信息与通信工程}
\end{cnabstract}
      % 中文摘要及关键词
%% -----------------------------------------------------------------
%% 北京邮电大学研究生学位论文 - 英文摘要
%% 格式要求（2025年1月15日规范）：
%% 1. “ABSTRACT”标题用 Times New Roman 字体三号字加粗，居中对齐，首行无缩进，段前0行，段后2行，单倍行距
%% 2. 摘要内容用四号 Times New Roman 字体，两端对齐，首行缩进2字符，段前0行，段后0行，固定值20磅
%% 3. “KEY WORDS:”四号字加粗，其后关键词为四号字，两端对齐，首行缩进2字符，段前2行，段后0行，固定值20磅
%% 4. 英文关键词与中文关键词对应，每一关键词之间用半角分号隔开（; ），最后一个关键词后不打标点符号
%% -----------------------------------------------------------------

\begin{enabstract}
The abstract should concisely present the essential contents and core methodology of the thesis. It highlights the research background, objectives, methodologies, significant theoretical contributions, experimental results, and conclusions.

This thesis focuses on key scientific challenges and engineering problems in next-generation information and communication networks. First, the existing literature and fundamental bottlenecks in the field are comprehensively summarized. Second, a novel adaptive optimization framework and distributed processing mechanism are proposed to achieve high robustness and low processing delay. Third, rigorous theoretical derivations are provided to characterize the convergence speed and capacity bounds of the proposed algorithms. Finally, extensive simulation studies and experimental evaluations on real-world datasets demonstrate the superior efficiency and reliability of the proposed scheme.

The results show that the proposed approach substantially reduces computational complexity while improving throughput and system accuracy, validating its effectiveness in realistic complex environments.

\enkeywords{Beijing University of Posts and Telecommunications; Graduate Degree Thesis; \LaTeX{} Template; Formatting Specification; Information and Communication Engineering}
\end{enabstract}
      % 英文摘要及关键词
\tableofcontents                     % 论文目录（自动排版至三级目录项）
%% -----------------------------------------------------------------
%% 北京邮电大学研究生学位论文 - 符号说明
%% 格式要求（2025年1月15日规范）：
%% 1. 标题用黑体三号字，居中对齐，首行无缩进，段前0行，段后2行，单倍行距
%% 2. 内容用宋体小四号字，段前和段后0行，固定值20磅
%% 3. 表明论文中使用符号的意义及单位（或量纲），一律采用国际通用标准
%% -----------------------------------------------------------------

\begin{nomenclature}
\begin{center}
  \renewcommand{\arraystretch}{1.5}
  \songti\zihao{-4}
  \begin{tabular}{>{\centering\arraybackslash}m{3.5cm} >{\centering\arraybackslash}m{6.5cm} >{\centering\arraybackslash}m{3.5cm}}
    \toprule
    符\quad 号 & 代\quad 表\quad 意\quad 义 & 单\quad 位 \\
    \midrule
    $A$       & 截面积；散热面积              & $\mathrm{m^2}$ \\
    $F$       & 力                          & $\mathrm{N}$ \\
    $e$       & 电子电荷量                  & $\mathrm{C}$ \\
    $I, i$    & 电流                        & $\mathrm{A}$ \\
    $m$       & 质量                        & $\mathrm{kg}$ \\
    $Q$       & 热量或流量                  & $\mathrm{kJ}$ \\
    $H$       & 焓值                        & $\mathrm{kJ/kg}$ \\
    $S$       & 熵值                        & $\mathrm{kJ/kg}$ \\
    $\lambda$ & 波长                        & $\mathrm{nm}$ \\
    $\sigma$  & 方差；电导率                & $\mathrm{S/m}$ \\
    $\mathbf{W}$ & 权重矩阵                 & -- \\
    $\mathcal{L}$ & 损失函数值              & -- \\
    \bottomrule
  \end{tabular}
\end{center}
\end{nomenclature}
    % 符号说明（三列表格）

%% ---------------- 正文部分（阿拉伯数字页码，从第1页重新编排） ----------------
\mainmatter
%% -----------------------------------------------------------------
%% 北京邮电大学研究生学位论文 - 第一章 绪论
%% 包含全部规范元素示例：多级标题、公式、图表、三线表、分图、长表续表、算法、脚注与文献引用
%% -----------------------------------------------------------------

\chapter{绪论}

\section{研究背景与意义}

正文文本采用宋体小四号，英文字母与阿拉伯数字采用 Times New Roman 字体，两端对齐，首行缩进 2 字符，段前 0 行，段后 0 行，行间距为固定值 20 磅\footnote{这是一个脚注示例。规范要求：脚注采用小五号宋体，西文采用 Times New Roman，每页重新按带圈数字序号编号。}。

随着新一代移动通信技术与人工智能的深度融合，大规模异构网络面临着计算资源受限、拓扑结构动态变化以及服务质量差异化等多重严峻挑战。针对这些问题，学术界和工业界开展了广泛的研究。文献~\cite{yuanxunlai2012}探讨了复杂系统演化机制；文献~\cite{vaswani2017attention}提出了基于自注意力机制的深度网络模型；相关国家与行业标准也在逐步完善~\cite{gbt7714, gbt25100}。

\subsection{国内外研究现状}

在国内外学者的共同努力下，网络协同计算与资源调度技术取得了长足的进步。根据学者钱学森等在系统学方面的经典论述~\cite{qianxuesen2001}，复杂系统的多维度建模需要综合考虑输入约束与多目标均衡。

\subsubsection{集中式调度机制的瓶颈}

集中式控制架构在小规模网络中表现良好，但在高并发场景下容易产生单点性能瓶颈与通信拥塞。

\paragraph{异构计算开销约束}
在边缘异构计算环境中，不同节点的处理能力差异显著，这要求系统必须具备自适应负载均衡能力。四级标题采用宋体小四号加粗（或宋体小四号），首行缩进 2 字符，题序与标题之间空 1 格。

\section{数学公式排版规范}

论文中公式应另起一行居中横排，公式编号按章连续编排，右端对齐，置于圆括号内，如式~(\ref{eq:system_model})与式~(\ref{eq:loss_func})。

设第 $k$ 个网络节点的信噪比为 $\gamma_k$，则该节点的信道容量可表示为香农公式形式：
\begin{equation}
  C_k = B \log_2 \left( 1 + \frac{P_k |h_k|^2}{\sigma^2 + \sum_{j \neq k} P_j |h_j|^2} \right)
  \label{eq:system_model}
\end{equation}
其中，$B$ 表示信道带宽，$P_k$ 为节点发射功率，$h_k$ 表示信道增益衰减系数，$\sigma^2$ 表示加性高斯白噪声功率。

系统总损失函数定义为预测误差与正则化项之和：
\begin{equation}
  \mathcal{L}(\mathbf{W}) = \frac{1}{N}\sum_{i=1}^N \ell(y_i, f(\mathbf{x}_i; \mathbf{W})) + \frac{\lambda}{2} \|\mathbf{W}\|_F^2
  \label{eq:loss_func}
\end{equation}

\section{图与子图排版规范}

规范要求：图应具有“自明性”，图的编号由“图”和从“1”开始的阿拉伯数字组成，如“图1-1”。图题置于图下方居中，中文字体为楷体五号，西文字体为 Times New Roman，固定值 20 磅。图序与图名之间空 1 格，图名中不允许使用标点符号，图名后不加标点。如有英文标题，使用 \texttt{\textbackslash bicaption} 宏命令输出。

\begin{figure}[htbp]
  \centering
  \includegraphics[width=0.28\textwidth]{resources/logo.png}
  \bicaption{北京邮电大学校徽标识图}{Beijing University of Posts and Telecommunications Emblem}
  \label{fig:bupt_logo}
  \fignote{注：校徽标识矢量图文件存放于 \texttt{resources/} 目录中。如有需要可对图片进行注释说明，可省略。}
\end{figure}

如果一个图由两个或两个以上分图组成时，各分图分别以 (a)、(b)、(c)…… 作为分图序，并须有分图名，分图名直接列在各自分图的正下方，主图名列在所有分图的下方正中，如图~\ref{fig:subfig_demo}所示。

\begin{figure}[htbp]
  \centering
  \begin{subfigure}[b]{0.35\textwidth}
    \centering
    \includegraphics[width=0.7\textwidth]{resources/logo.png}
    \caption{子系统节点模型 A}
    \label{fig:sub_a}
  \end{subfigure}
  \hspace{1.5cm}
  \begin{subfigure}[b]{0.35\textwidth}
    \centering
    \includegraphics[width=0.7\textwidth]{resources/logo.png}
    \caption{子系统节点模型 B}
    \label{fig:sub_b}
  \end{subfigure}
  \bicaption{多节点协同系统拓扑对比图}{Comparison of Multi-Node Collaborative System Topologies}
  \label{fig:subfig_demo}
\end{figure}

\section{表格与三线表排版规范}

表的编排一般采用国际通行的三线表。表编号由“表”和从“1”开始的阿拉伯数字组成，如“表1-1”。表题置于表上方居中，使用楷体五号，固定值 20 磅。表序与表名之间空 1 格，表名后不加标点。

\begin{table}[htbp]
  \centering
  \bicaption{系统运行核心参数设置表}{Core Parameter Configurations of the Evaluated System}
  \label{tab:system_params}
  \renewcommand{\arraystretch}{1.3}
  \begin{tabular}{lccc}
    \toprule
    参数名称 & 符号表示 & 默认基准值 & 取值范围 \\
    \midrule
    载波中心频率 & $f_c$ & $3.5\ \mathrm{GHz}$ & $[2.0, 6.0]$ \\
    系统带宽 & $B$ & $100\ \mathrm{MHz}$ & $[20, 200]$ \\
    基站发射功率 & $P_t$ & $43\ \mathrm{dBm}$ & $[30, 46]$ \\
    噪声功率谱密度 & $N_0$ & $-174\ \mathrm{dBm/Hz}$ & 常数 \\
    节点移动速度 & $v$ & $5\ \mathrm{m/s}$ & $[0, 30]$ \\
    \bottomrule
  \end{tabular}
  \tablenote{数据来源：如需对数据来源进行说明，请参照此格式排版，首行缩进 2 字符。}
\end{table}

\section{跨页长表格与续表排版}

如某个表需要转页接排，在随后各页上应重复表的编号与表头，如“续表1-2 表名”，英文则用“Table 1-2 Table name (continued)”，且续表均应重复表头。本模板通过 \texttt{longtable} 宏包原生支持该规范要求，如表~\ref{tab:longtable_example}所示。

\begin{center}
\renewcommand{\arraystretch}{1.3}
\begin{longtable}{>{\centering\arraybackslash}m{2.5cm} >{\centering\arraybackslash}m{3.5cm} >{\centering\arraybackslash}m{3.5cm} >{\centering\arraybackslash}m{4.0cm}}
  \caption{大规模测试样本属性统计长表} \label{tab:longtable_example} \\
  \toprule
  样本编号 & 特征维度一 & 特征维度二 & 标注分类结果 \\
  \midrule
  \endfirsthead

  \multicolumn{4}{c}{{\songti\zihao{5}\kaishu 续表1-2 大规模测试样本属性统计长表}} \\
  \toprule
  样本编号 & 特征维度一 & 特征维度二 & 标注分类结果 \\
  \midrule
  \endhead

  \bottomrule
  \endfoot

  \bottomrule
  \endlastfoot

  Sample-001 & $12.45 \pm 0.12$ & $0.982$ & 正常传输状态 \\
  Sample-002 & $15.89 \pm 0.25$ & $0.941$ & 正常传输状态 \\
  Sample-003 & $28.34 \pm 0.88$ & $0.612$ & 弱信道衰落状态 \\
  Sample-004 & $31.20 \pm 1.05$ & $0.485$ & 强干扰阻塞状态 \\
  Sample-005 & $14.12 \pm 0.18$ & $0.963$ & 正常传输状态 \\
  Sample-006 & $16.78 \pm 0.32$ & $0.925$ & 正常传输状态 \\
\end{longtable}
\end{center}

\section{算法伪代码排版规范}

算法环境在论文中通常按章编号，如“算法1-1”。下面展示了一个规范的自适应梯度优化算法编写示例。

\begin{table}[htbp]
  \centering
  \begin{minipage}{0.9\textwidth}
    \rule{\linewidth}{1.0pt}\vspace{3pt}
    \noindent\textbf{算法 1-1} \quad 动态权重自适应迭代优化算法 (AdaOpt)\\[-6pt]
    \rule{\linewidth}{0.5pt}\vspace{4pt}
    \textbf{输入：} 训练样本集合 $\mathcal{D}$，学习率 $\alpha$，衰减系数 $\beta$，最大迭代步数 $T$；\\
    \textbf{输出：} 优化收敛权重参数 $\mathbf{W}^*$。
    \begin{enumerate}[leftmargin=2.0em, itemsep=2pt, topsep=2pt]
      \item[1:] 初始化模型参数 $\mathbf{W}_0 \leftarrow \mathbf{0}$，一阶动量 $\mathbf{m}_0 \leftarrow \mathbf{0}$，时间步 $t \leftarrow 0$；
      \item[2:] \textbf{while} $t < T$ 且 $\|\nabla \mathcal{L}(\mathbf{W}_t)\| > \epsilon$ \textbf{do}
      \item[3:] \quad $t \leftarrow t + 1$；
      \item[4:] \quad 从数据集中抽取当前微批次样本 $\mathcal{B}_t \subset \mathcal{D}$；
      \item[5:] \quad 计算损失函数关于当前权重的梯度值 $\mathbf{g}_t \leftarrow \nabla_{\mathbf{W}} \mathcal{L}(\mathbf{W}_{t-1}; \mathcal{B}_t)$；
      \item[6:] \quad 更新指数移动平均动量 $\mathbf{m}_t \leftarrow \beta \mathbf{m}_{t-1} + (1 - \beta) \mathbf{g}_t$；
      \item[7:] \quad 更新网络权重参数 $\mathbf{W}_t \leftarrow \mathbf{W}_{t-1} - \alpha \mathbf{m}_t$；
      \item[8:] \textbf{end while}
      \item[9:] \textbf{return} 最终模型权重 $\mathbf{W}^* = \mathbf{W}_t$。
    \end{enumerate}
    \rule{\linewidth}{1.0pt}
  \end{minipage}
\end{table}

\section{本章小结}

本章详细阐述了论文的研究背景、理论意义与工程应用价值，并给出了北邮研究生学位论文在各级标题、数学公式、双语图表、三线表、长表续表、算法伪代码及参考文献引用方面的规范排版标准。后续章节将围绕核心创新点展开深入讨论。
    % 第一章 绪论
%% -----------------------------------------------------------------
%% 北京邮电大学研究生学位论文 - 第二章 理论与方法
%% -----------------------------------------------------------------

\chapter{理论分析与系统建模}

\section{引言}

本章重点阐述论文所提方法的核心理论架构与系统数学建模。在第一章综述的基础上，针对现有算法在复杂信道条件下的性能退化现象，提出了一种多尺度协同优化框架。

\section{系统模型与问题建模}

考虑由 $M$ 个接入节点与 $K$ 个边缘终端组成的分布式处理网络。节点间通过无线正交频分复用链路进行互联互通。系统的整体效用函数表示为：
\begin{equation}
  U(\mathbf{p}, \mathbf{w}) = \sum_{k=1}^K w_k \log \left( 1 + \gamma_k(\mathbf{p}) \right) - \sum_{m=1}^M \alpha_m p_m
  \label{eq:utility}
\end{equation}
其中，$\mathbf{p} = [p_1, \dots, p_M]^T$ 为各节点发射功率向量，$\mathbf{w}$ 为服务优先级加权向量，$\alpha_m$ 为功耗代价因子。

\subsection{凸松弛与对偶分解}

由于式~(\ref{eq:utility})中存在非凸耦合约束，直接求解属于 NP-hard 难题。本文采用对偶分解法将原始全局优化问题转化为一系列局部子问题：
\begin{equation}
  \mathcal{D}(\boldsymbol{\lambda}) = \max_{\mathbf{p} \in \mathcal{P}} \left\{ U(\mathbf{p}, \mathbf{w}) - \boldsymbol{\lambda}^T (\mathbf{A}\mathbf{p} - \mathbf{b}) \right\}
  \label{eq:dual}
\end{equation}

\section{本章小结}

本章建立了多节点分布式系统的效用最优化理论模型，证明了对偶间隙的存在条件与收敛性边界，为后续算法设计与实证分析提供了理论依据。
          % 第二章 理论分析与系统建模
%% -----------------------------------------------------------------
%% 北京邮电大学研究生学位论文 - 第三章 实验与结果分析
%% -----------------------------------------------------------------

\chapter{实验验证与结果分析}

\section{实验环境与数据集}

为了验证本文所提理论模型与算法方案的优越性，本章在仿真测试平台与实际工程语料上进行了充分的实验评估。基准测试数据集包含来自真实通信链路的连续采样样本。

\begin{table}[htbp]
  \centering
  \bicaption{不同算法在测试集上的性能对比指标}{Performance Comparison of Different Algorithms on the Benchmark Dataset}
  \label{tab:exp_results}
  \renewcommand{\arraystretch}{1.3}
  \begin{tabular}{lcccc}
    \toprule
    算法方案 & 运行时间 ($\mathrm{ms}$) & 吞吐量 ($\mathrm{Mbps}$) & 准确率 (\%) & 能效比 ($\mathrm{bit/J}$) \\
    \midrule
    基线算法 A~\cite{nalamura2014} & 45.2 & 124.6 & 88.5 & 312.4 \\
    基线算法 B~\cite{mahuan2011}   & 38.7 & 142.1 & 91.2 & 389.0 \\
    本文所提方案                    & \textbf{21.4} & \textbf{198.5} & \textbf{96.8} & \textbf{542.7} \\
    \bottomrule
  \end{tabular}
  \tablenote{注：表中粗体数值表示在当前对比评估项目下的最优性能指标。}
\end{table}

\section{性能指标对比与讨论}

由表~\ref{tab:exp_results}可以看出，相较于传统对比方案，本文算法在运算速度上提升了约 44\%，吞吐量提升了 39\% 以上。这一结果充分体现了自适应松弛机制在处理异构数据流时的明显优势。

\section{本章小结}

本章通过定性与定量对比实验，全面评估了新方法在不同场景下的综合表现，各项指标均达到了预期设计要求。
     % 第三章 实验验证与结果分析
%% -----------------------------------------------------------------
%% 北京邮电大学研究生学位论文 - 第四章 总结与展望
%% -----------------------------------------------------------------

\chapter{总结与展望}

\section{全文工作总结}

本文围绕研究生学位论文的研究课题，展开了系统深入的理论探索、机制构建与工程实证。论文的主要研究成果归纳如下：
\begin{enumerate}[leftmargin=2.5em, itemsep=4pt, label=（\arabic*）]
  \item 深入剖析了复杂信息处理系统的瓶颈机理，建立了端到端效用优化数学模型；
  \item 提出了一种高效低开销的自适应迭代求解机制，给出了理论收敛性证明；
  \item 搭建了原型评测系统并在多组真实场景下开展了对比实验，验证了算法的先进性与实用性。
\end{enumerate}

\section{未来研究展望}

尽管本文取得了一定阶段性成果，但在以下方面仍具有进一步拓展与深化的空间：
\begin{enumerate}[leftmargin=2.5em, itemsep=4pt, label=（\arabic*）]
  \item 进一步探索在大规模非平稳动态拓扑环境下的自愈合容错能力；
  \item 将所提理论框架推广应用至跨域异构工业互联网及天地一体化信息网络中。
\end{enumerate}
      % 第四章 总结与展望

%% ---------------- 参考文献部分（GB/T 7714-2015 格式） ----------------
\bibliographystyle{gbt7714-numerical}
\bibliography{Bib/thesis}

%% ---------------- 附录与后置章节 ----------------
%% -----------------------------------------------------------------
%% 北京邮电大学研究生学位论文 - 附录
%% 格式要求（2025年1月15日规范）：
%% 1. 标题用黑体三号字，居中对齐，首行无缩进，段前0行，段后2行，单倍行距
%% 2. 附录内容编于正文后，公式、图表重新按字母编号（如 式(A-1)、图A-1 等）
%% -----------------------------------------------------------------

\appendix

\chapter{数学推导与辅助定理证明}

\section{定理 1 的收敛性证明}

此处补充正文中过于冗长的公式推导与引理证明。

设序列 $\{x_k\}$ 满足以下递归不等式：
\begin{equation}
  \|x_{k+1} - x^*\|^2 \le (1 - \rho)\|x_k - x^*\|^2 + \epsilon_k
  \label{eq:appendix_proof}
\end{equation}
其中 $0 < \rho < 1$，误差项满足 $\sum_{k=1}^\infty \epsilon_k < \infty$。根据超鞅收敛定理，可知序列 $\{x_k\}$ 以概率 1 强收敛于最优点 $x^*$。

\section{重要辅助参数表}

\begin{table}[htbp]
  \centering
  \bicaption{仿真实验补充配置参数表}{Supplementary Configuration Parameters for Simulation}
  \label{tab:appendix_table}
  \renewcommand{\arraystretch}{1.3}
  \begin{tabular}{lll}
    \toprule
    参数名称 & 配置值 & 单位说明 \\
    \midrule
    蒙特卡洛仿真次数 & 10000 & 次 \\
    信道衰落模型 & 瑞利平坦衰落 & -- \\
    天线增益 & 15 & $\mathrm{dBi}$ \\
    \bottomrule
  \end{tabular}
\end{table}
        % 附录
%% -----------------------------------------------------------------
%% 北京邮电大学研究生学位论文 - 致谢
%% 格式要求（2025年1月15日规范）：
%% 1. 标题用黑体三号字，居中对齐，首行无缩进，段前0行，段后2行，单倍行距，中间空1个全角空格
%% 2. 内容用宋体小四号，两端对齐，首行缩进2字符，段前和段后0行，固定值20磅
%% 3. 对给予各类资助、指导和协助完成研究工作以及提供各种对论文工作有利条件的单位和个人表示感谢。致谢应实事求是，切忌浮夸之词。
%% -----------------------------------------------------------------

\begin{acknowledgement}
在研究生阶段的学习与学位论文撰写过程中，我得到了来自导师、实验室同门、家人以及各方师友的悉心指导与全力支持，在此谨向所有给予我关怀与帮助的人表达最诚挚的谢意。

首先，衷心感谢我的指导教师在学术方向上的悉心指引与严谨教诲。从论文选题、方案构思到具体实验与论文定稿，导师倾注了大量的心血，其求真务实、严谨治学的学术风范令我终身受益。

感谢课题组全体老师和同学在日常科研与实验工作中的支持与合作。在浓厚热烈的学术氛围中，大家互相启发、共同进步，度过了充实而难忘的科研岁月。

感谢北京邮电大学为我们提供了良好的科研平台、丰富的文献资源与优良的学习环境。

最后，深深地感谢我的家人一直以来的默默付出与坚定支持，他们的理解与关爱是我克服困难、不断前行的最大动力。
\end{acknowledgement}
% 致谢
%% -----------------------------------------------------------------
%% 北京邮电大学研究生学位论文 - 攻读学位期间取得的创新成果目录
%% 格式要求（2025年1月15日规范）：
%% 1. 标题为“攻读学位期间取得的创新成果目录”（注：最新规范不带双引号）
%% 2. 标题黑体三号字，居中对齐，首行无缩进，段前0行，段后2行，单倍行距
%% 3. 内容用宋体小四号，固定值20磅
%% 4. 按取得成果的时间顺序，列出本人在攻读学位期间取得的创新成果清单
%% -----------------------------------------------------------------

\begin{publication}

\section*{一、学术期刊与会议论文}
\begin{enumerate}[leftmargin=2.5em, itemsep=4pt, label={[\arabic*]}]
  \item \textbf{张三}, 李四, 王五. 基于自适应松弛机制的高效协同计算方法研究[J]. 软件学报, 2024, 35(8): 3610--3625. (CCF-A类中文核心期刊)
  \item \textbf{San Zhang}, Si Li, Wu Wang. Adaptive optimization framework for multi-tier edge networks[C]// Proceedings of the IEEE International Conference on Computer Communications (INFOCOM). IEEE, 2024: 1120--1129. (CCF-A类国际顶级会议)
  \item Si Li, \textbf{San Zhang}, Wu Wang. Robust distributed learning over fading channels[J]. IEEE Transactions on Mobile Computing, 2023, 22(11): 6580--6594. (CCF-A类国际期刊)
\end{enumerate}

\section*{二、申请与授权专利}
\begin{enumerate}[leftmargin=2.5em, itemsep=4pt, label={[\arabic*]}]
  \item 李四, \textbf{张三}, 王五. 一种面向边缘计算的动态自适应资源调度方法及系统: CN202310888888.8[P]. 2024-03-15. (实质审查阶段)
  \item \textbf{张三}, 李四. 一种高可靠异构网络数据传输优化方法: CN202210666666.6[P]. 2023-11-20. (已授权)
\end{enumerate}

\section*{三、参与科研项目}
\begin{enumerate}[leftmargin=2.5em, itemsep=4pt, label={[\arabic*]}]
  \item 国家自然科学基金重点项目：“面向下一代移动通信的高性能协同计算理论与关键技术”（项目编号：62231000），2022.01--2026.12，主要学生骨干参与。
  \item 国家重点研发计划课题：“异构网络自适应调度平台研制与应用示范”（项目编号：2023YFB2900000），2023.01--2025.12，主要参与人员。
\end{enumerate}

\end{publication}
    % 攻读学位期间取得的创新成果目录
%% -----------------------------------------------------------------
%% 北京邮电大学研究生学位论文 - 作者学习经历
%% 格式要求（2025年1月15日规范）：
%% 1. 标题为“作者学习经历”，黑体三号字，居中对齐，首行无缩进，段前0行，段后2行，单倍行距
%% 2. 内容用宋体小四号，两端对齐，固定值20磅
%% 3. 按时间顺序，从高中阶段开始填写
%% -----------------------------------------------------------------

\begin{experience}

\begin{center}
  \renewcommand{\arraystretch}{1.5}
  \songti\zihao{-4}
  \begin{tabular}{|>{\centering\arraybackslash}m{3.5cm}|>{\centering\arraybackslash}m{2.0cm}|>{\centering\arraybackslash}m{4.5cm}|>{\centering\arraybackslash}m{4.0cm}|}
    \hline
    起止年月 & 类\quad 别 & 学习单位名称 & 学科/专业名称 \\
    \hline
    2014.09 -- 2017.06 & 高中 & 北京市第一中学 & -- \\
    \hline
    2017.09 -- 2021.06 & 本科 & 北京邮电大学 & 通信工程 \\
    \hline
    2021.09 -- 2024.06 & 硕士 & 北京邮电大学 & 信息与通信工程 \\
    \hline
    2024.09 -- 至今 & 博士 & 北京邮电大学 & 智能科学与技术 \\
    \hline
  \end{tabular}
\end{center}

\end{experience}
      % 作者学习经历

\end{document}
